% Nouveau EGA V (cote 25): archive cote 25, site page 174 (PDF page 175).
% typescript. Draft compared: none.
% AI-transcribed by Claude Opus 5.5 (claude-opus-5-5, adaptive thinking, effort high), Message Batches API, batch ega25b (msgbatch_01A7n5YuXubUPNK1zsepi8RG), 2026-11-13; not yet checked by eye.
\scanpage{25/174}
\margin{\unsure{IV|1/U}} \margin{nouveau EGA V} \mspage{1}

17.16. \uline{Formes quadratiques lisses}.

16.1.Soit $Q$ \unsure{une} section de $\underline{Sym}^2(\underline{E})$, ($\underline{E}$ loc libre de rang fini),
de sorte que Q \xout{\unsure{définit}} (qu'on peut aussi interpréter comme une forme
quadratique sur $\overset{\hand{$\vee$}}{\underline{E}}$ ) \xout{\illegible{3} \unsure{section}} définit une forme biliné-%
aire symétrique B sur $\overset{\hand{$\vee$}}{\underline{E}}$, d'où un homomorphisme $\overset{\hand{$\vee$}}{\underline{E}} \longrightarrow \underline{E}$ , et passant aux
Modules déterminants on trouve $\text{dét}\ \overset{\hand{$\vee$}}{\underline{E}} \longrightarrow \det \underline{E}$ , d'où enfin une section
\[ d(Q) \in \Gamma(\ \text{dét}\ (E)^{\otimes 2}) \]
appelée (sauf erreur) \uline{discriminant de la forme quadratique} Q. Lorsque
$\underline{E} = \underline{O}_S^{\ \hand{$n$}}$ , c'est simplement la section de $\underline{O}_S$ déterminant de la matrice
à n lignes et n colonnes, à coeff dans $\Gamma\underline{O}_S$ , qui exprime Q . Notons tout
de suite:

\uline{Proposition} 16.2. Supposons que pour tout $s \in S$, on ait car $k(s) \neq 2$ \uline{ou}
rang. $E\otimes k(s)$ pair. Alors, pour que Q soit lisse, il f et s que Q soit
``non dégénérée'' i.e. que d(Q) soit inversible. \hand{(}D'ailleurs la condition est
suffisante sans restriction sur S ni $\underline{E}$\hand{)}.

La dernière assertion n'est mise que pour la commodité des référen-%
ces éventuelles, car elle est trivialement contenue dans la première,
l'hypothèse que Q est non dégénérée implique en effet que $\underline{E}$ est de rang
pair en tout point dont la car est 2. D'ailleurs, elle se démontre direc-%
tement sans considération de caractéristique, car (se plaçant sur un corps
\inserted{\unsure{alg clos}} et prenant une base) elle exprime que les $\frac{\hand{$\partial$} Q}{\hand{$\partial$} X_i}$ n'ont pas de zéro\xout{s} commun\xout{s}
non trivial, à fortiori \unsure{ils} n'ont pas de zéro commun non trivial avec Q.
Si car $k \neq 2$, la réciproque est vraie, car en vertu de la formule
\[ 2Q = \hand{$\sum$} X_i \frac{\hand{$\partial$} Q}{\hand{$\partial$} X_i} \]
on voit alors que tout zéro commun aux $\frac{\hand{$\partial$} \unsure{Q|E}}{\hand{$\partial$} X_i}$ est aussi un zéro de Q. Enfin,
si car $k = 2$ , \xout{si E est de rang pair} alors la forme bilinéaire B

\margin{174}
